\section{Portability Across Metrics}\label{sec:results-cross-metric}
Q3 asks whether a preferred-policy conclusion survives a change of objective.
This section uses the saved post-campaign cross-metric extension and new,
explicitly post-hoc M2 sensitivities on its frozen inputs. No serving
execution is added. All 55 unordered pairs among the 11 eligible metrics
are evaluated for each of three sources and six regions: 990 conditions.
These conditions share outcomes and are not independent observations.

\subsection{Historical correlations and classification counts}
The historical median Kendall tau-b is 0.419, with range $[-0.506,1.000]$;
202 of 990 conditions have negative correlation. The deterministic top
choice agrees in 316 conditions (31.9\%). Thus the objective can change
both the full ordering and the selected policy ID. The median describes
this enumerated set of metric pairs; it is not an expectation over a
population of future deployment objectives.

The 54,450 policy-pair records divide into 45,800 same-order cases,
5,205 microscopic sign changes, 201 practical sign changes with a wide
saved interval, and 3,244 saved-supported practical disagreements. The
historical screened BH procedure flags 3,207 of the last group, removing
37. These are preserved procedural counts. M2 withdraws the stronger
interpretation that they are statistically established practical effects
at a controlled error rate. The limitations of the component tail
quantities, selection and dependence apply here as in Q2.

No condition-level ranking is undefined, but that does not imply complete
policy/window observations. Nor do 54,450 records supply 54,450 samples:
each window contributes to multiple policies, metrics and regions. The
number of pairwise records is useful for describing the evaluation grid,
not for justifying precision of a claim about general serving workloads.

\subsection{Tied optima versus deterministic policy selection}
Let $T_m=\{p:\bar y_{p,m}=\max_j\bar y_{j,m}\}$ after direction
normalization. The historical rule selects the smallest policy ID in
$T_m$. A set-based question instead asks whether $T_a\cap T_b$ is
nonempty: is there any policy maximizing both observed objectives?
Equality of the two sets asks a stricter question. Neither set criterion
selects a policy without an additional preference rule.

Among 1,980 metric sides, 1,010 have tied maxima. In the M2 post-hoc
analysis, the maximizing sets intersect in 572/990 conditions (57.8\%),
and are identical in 179/990 (18.1\%). Only 257 conditions have a unique
maximizer on both sides. Mean Jaccard overlap,
$|T_a\cap T_b|/|T_a\cup T_b|$, is 0.329. These statistics use exact
aggregate equality, matching the historical arithmetic; no tolerance was
tuned to improve agreement. Near-optimal sets under a deployment-specific
utility tolerance are a different, unmeasured object.

\begin{table}[t]\centering
\caption{Historical and post-hoc agreement summaries on the same 990
cross-metric conditions. Different definitions answer different questions.}
\label{tab:top-sets}
\begin{tabular}{p{0.58\linewidth}rr}\toprule
Agreement definition & Count & Fraction\\\midrule
Historical deterministic ID & 316 & 31.9\%\\
Post-hoc nonempty maximizing-set intersection & 572 & 57.8\%\\
Post-hoc equality of maximizing sets & 179 & 18.1\%\\
Both objectives have unique maxima (eligibility) & 257 & 26.0\%\\\bottomrule
\end{tabular}\end{table}
The additional 256 intersecting-set conditions show why the complement
of 31.9\% cannot be called disagreement between unique optima. A common
best policy may exist even when the ID rule picks different members of
unequal sets. Conversely, nonempty intersection can be permissive when
one objective ties many policies. Reporting intersection alongside set
size and equality avoids presenting 57.8\% as a corrected version of
31.9\%. Both are valid descriptions of different selection questions.

\subsection{Matched support and the meaning of a metric difference}
Historical policy means omit undefined observations separately. To compare
$a$ and $b$, the difference of those means is
\begin{equation}
 d_{\mathrm{sep}}=
 \frac{\sum_{w\in\mathcal W_a}y_{w,a}}{|\mathcal W_a|}
 -\frac{\sum_{w\in\mathcal W_b}y_{w,b}}{|\mathcal W_b|},
\end{equation}
whereas the bootstrap difference averages $y_{w,a}-y_{w,b}$ over
$\mathcal W_a\cap\mathcal W_b$. Without equal support, its interval is
not an interval for $d_{\mathrm{sep}}$. This is an estimand mismatch,
not simply a wider or narrower error bar. M2 therefore avoids assigning
that interval to the separately aggregated point difference.

We perform two distinct post-hoc checks. First, each metric-pair ranking
is recomputed over windows where all 11 policies have both metrics.
This common-panel support retains 33--40 windows per condition. The
median tau becomes 0.41849, 202 correlations remain negative, and the ID
agreement remains 316/990. Individual tau values change in 81 conditions.
The historical median remains reported as 0.419; the new value is a
sensitivity estimate with a different eligibility rule.

Second, for every policy-pair record, each metric's policy means are
recomputed over that pair's observed-window intersection. This retains
more observations than requiring the entire panel. Differences change
in 5,890 of 54,450 records, including 2,720 of the 3,244 saved-supported
records. All 3,244 retain opposite signs and both 10\% point-margin gates;
for the affected saved-supported records, the matched signs also agree
with the saved interval signs. This supports robustness of those
\emph{descriptive gates} in this corpus. It neither validates their
$p$-values nor supplies fresh intervals for newly eligible cases. We do
not report a replacement total of statistically supported disagreements.

Matching support also changes the scientific target. A latency comparison
restricted to windows on which both policies finish requests does not
represent windows where one finishes none. Within an observed window,
each policy's latency may still describe a different set of completed
requests. Matched windows therefore do not eliminate request-level
survivorship selection. ANWG answers an arrival-based question; conditional
latency answers a completion-based question. Neither should be silently
substituted for the other.

\subsection{Implications for objective-dependent selection}
A policy can have low latency among its completed requests while providing
little arrival-normalized deadline utility. Throughput and TTFT measure
other aspects of service; priority weighting changes whose outcomes count
most. Some disagreement between such rankings is consequently compatible
with internally consistent measurements. Negative tau does not by itself
identify a measurement error or a universally inferior scheduler.

For evaluation practice, the objective, denominator and eligibility rule
must accompany the selected policy. If the intended decision is to attain
arrival deadlines, completion-conditioned speed is insufficient evidence.
If several policies attain the same measured objective, the benchmark can
report that set rather than implying that an arbitrary ID is uniquely
best. A deployment may use another objective to break ties, but this
corpus does not validate that decision rule, its cost or its portability.

The saved SLO extension in Section~\ref{sec:results-slo} changes the
benchmark's synthetic deadline convention rather than merely relabeling
an unchanged latency statistic. Its 13/30 top-choice changes and six
reversal-direction changes among 72 nonidentity checks show dependence on
that convention within its selected scope. The extension excludes LOW
because its design targeted saved reversal cases; no conclusion about
LOW's deadline sensitivity follows. These distinctions answer Q3 without
turning retrospective sensitivity into a prospective deployment guarantee.
